% New GPT V 5.4 in action. It is very slow but hopefully provides a better rendering.

% TODO: Settle "what is least is faithful also in much" vs "small things and great things"

\section{Ninth Sunday after Trinity: Faithful in Little and in Much}



\begin{quote}

Luke 16:10–17. He who is faithful in little is faithful also in much, and he who is unrighteous in little is unrighteous also in much. If then you have not been faithful with unrighteous Mammon, who will entrust to you the true riches? And if you have not been faithful in what belongs to another, who will give you what is your own? No servant can serve two masters; for either he will hate the one and love the other, or else he will cling to the one and despise the other. You cannot serve God and Mammon.
But the Pharisees also heard all this, and they were lovers of money, and they mocked him. And he said to them, “You are those who justify yourselves before men, but God knows your hearts; for what is exalted among men is an abomination before God. The Law and the Prophets were until John; since that time the kingdom of God is proclaimed through the gospel, and everyone presses into it by force. But it is easier for heaven and earth to pass away than for one tittle of the Law to fail.”

\end{quote}

\bigskip


True Christianity calls for true faithfulness, since we are God’s servants and God’s stewards. We must give account to the Lord for all things; therefore there must be faithfulness in little and in much, in temporal things and in spiritual things.

But all our faithfulness must be toward the Lord. It is of no use to be a faithful servant of Mammon, toiling with the spirit of a true slave and with unbroken faithfulness under Mammon’s yoke all the days of one’s life, without ever lifting a finger to burst those bonds. If any man would be a Christian, his faithfulness must be toward the Lord, both in temporal things and in spiritual things. It will not do, as so many seemingly godly people do, to serve Mammon in temporal things and the Lord in spiritual things; in the end it is all nothing but the service of Mammon, all of it, with a hypocritical show of piety and spiritual trifles without power and truth.

There are states of soul more dangerous than that in which a man tries to serve the Lord in spiritual things, and yet remains in Mammon’s service in earthly things. That was the Pharisees’ misfortune. They would serve the Lord without being truly set free from the world with their whole heart; then their life became outwardly honorable and righteous, and so they were “not like other men, robbers, unjust, adulterers;” but then in covetousness they took back what they had forsaken in other ways.

So it is in our own day as well with the children of bondage, with those who, terrified by the Law and by wrath, are driven to an outward renunciation of the world and its pleasures, without their hearts having been set free by grace. Hard and strict, bitter and loveless, they deny themselves all outward joys of the world; but love of money keeps their hearts bound in the slavery of Mammon; and while they would serve God in spiritual things, and yet are not set free from the world in earthly things, they become in reality nothing but Mammon’s slaves, despite all the show of godliness.

There is therefore no more dangerous vice than the love of money. It is so proper and respectable; it is so frugal and sensible; it looks like the most perfect renunciation of the world, it looks like unbroken faithfulness in the stewardship of earthly goods, but in reality it is a complete surrender of the heart to the world and to all its desires; for the miser’s heart secretly enjoys everything that his money can buy for him, without spending any of his precious treasure. He delights in the sight of the money, because he knows what it could buy; but he keeps the money, because he knows that once he has spent it, he can enjoy it no longer.

In this way, he who would be faithful to two masters is compelled to love the one and despise the other. For he who joins the fear of God with love of money, he mocks and despises the God whom he serves with his mouth and honors with his lips, while his heart clings to the world and to Mammon. Woe unto those in this double bondage; a slave of God and a slave of the world; driven by the Law to serve as slaves under God, bound by the world to serve under Mammon in the dreadful vice of love for money. And yet dreadfully many people are in exactly this condition. O that they would remember that “the slave does not remain in the house forever, but the Son remains in the house forever!” Such a slave leads a tormented and miserable life in bitter anguish, who would serve both God and the world, and at last it ends in eternal torment; for “the slave does not remain in the house forever.”

No, a man cannot serve two masters.

It is not possible to be faithful to Mammon in earthly things and to be faithful to the Lord in spiritual things. You must be unfaithful to Mammon and break his yoke and renounce his service in earthly things, if you would be the Lord’s servant in true faithfulness. You must become an unfaithful steward toward Mammon, if you are to become a faithful steward for the Lord.

How difficult this is! And yet it must be learned in Mammon’s school. Mammon says: be frugal, be careful, be diligent, make use of every opportunity to gain money, flee and shun every opportunity to spend it; remember that you may soon become ill, think that in some years you will become old and frail, therefore save and scrape together all you can. It sounds so reasonable; it rings so sensible; it looks like faithfulness; but alas, it is faithfulness toward Mammon and unfaithfulness toward God; whereas your calling, O Christian soul, should be faithfulness toward God and unfaithfulness toward Mammon.

How then shall I be faithful, you say, if thrift and carefulness are not faithfulness? Does that mean we should be wasteful, slothful, and neglectful? 

No, you know far better than that. 

That too is only a new way of slaving under your carnal lusts. It is no faithfulness toward the Lord. But come with me to Jesus’ cross and learn there first what it is to love: to love God and to love men; then you will begin to understand what it is to be unfaithful to Mammon and faithful to God. If grace makes you free, free from God’s wrath, free from the slavery of the world, then you can become faithful to the Lord in all things, faithful in earthly things, faithful in heavenly things, faithful in little and faithful in much. 

For faithfulness stands in this: to love and to use all things in the service of love.

Thus a child of God thoroughly outwits Mammon, and this is true unfaithfulness toward him, when a child of God is both industrious and thrifty, both frugal and prudent, makes use of a rightful opportunity to gain money and shuns an evil opportunity to squander it, in order to have enough to share with those who are in need. That is what it means to be unfaithful to Mammon and faithful to the Lord. For he says: “Whatever you have done for these little ones, you have done for me.”

If any man would be faithful in little, faithful to the Lord and unfaithful toward Mammon, he uses his earthly means in the service of love, in the service of the kingdom of God. If he is truly faithful in this, then he is also faithful in much, in faith, in hope, in love. And whoever does not know how to do good and to distribute what he has, does not understand either the great divine stewardship, whose main point is this, that the Lord has delight in mercy and not in sacrifice. It is of no use to say: Yes, I know well enough of God’s love, but I have nothing but crumbs for Lazarus and for my suffering fellow men and for the spread of God’s kingdom among Jews and Gentiles. That is not faithfulness toward the Lord; it is faithfulness toward Mammon.

Oh, that we might become a little more faithful in little! We are quick to imagine that we are so faithful in much that we have no need to be faithful in little. We think we have such deep insight into God’s revelation and the doctrine of his Word that we need not concern ourselves with charity and mission and all the things that require sacrifice and generosity. We serve the Lord in much, and then it may be all the same to us whether we are faithful in little. We are zealous in prayer meetings and assemblies; we are so faithful in these that home may be neglected. God have mercy on us! If we have become too great to be faithful in little, then the Lord has no regard for our imagined faithfulness in much. The Lord will say to us as to the Pharisees: “You are those who justify yourselves before men, but God knows your hearts; for what is exalted among men is an abomination before God.”

But blessed is he who is faithful in little; for in his own time the Lord shall say unto him: 

“Well done, good and faithful servant, 

you have been faithful over little, 

I will set you over much; 

enter into the joy of your Lord!”

